\chapter[The Left Hand]{Adding the Left Hand to Unlock the Guitar's Full Potential}

In this chapter, we will explore the various uses of the left hand to select notes on the fretboard. 

\section{The job of the left hand}

In contrast with the right hand, the job of the left hand is focused on note and chord selection. If the left hand was absent, we would only be able to play six different notes on the guitar, and all our music would sound painfully similar. A well-functioning left hand allows us to rapidly select from a vast repository of different notes, which in turn lets us produce beautiful melodies, electrifying scales and arpeggios, and sonorous chords, just to name a few possibilities. 

\section{Left hand posture and mechanics}

To prepare the left hand for proper usage, move the left elbow slightly to the left until the arm feels more open. Next, place the thumb on center of the guitar neck, directly behind the second fret box. The fingers should naturally find their way over the strings. Start by letting the palm of the left hand touch the strings, then arch the fingers so that only the fingertips are touching the strings. Once this becomes comfortable, try placing all four fingertips (index, middle, ring, and pinky) on the bottom nylon string, starting with index finger in the first fret box and letting the other three fingers fall naturally where they may. 

\fbox{\textbf{Tip:} Don't worry if you can't reach the fourth fret with the pinky.} 

%Image of the four fingers arched and resting on the 1st string E

\section{``Corner" of the fret box}

Before you try pushing the string down, you must learn that there is a specific place in each fret box that is ideal for placing the finger. The easiest place to push the string down is located as close to the metal fret bar as possible without going over it. This area requires the least amount of force to get the string down. It is harder to push down the string if you touch the middle, or worse, the left side of the fret box. Remember this special area as the ``corner" of the fret box, and always place your fingers there before pushing down. 

%Two images showing correct and incorrect finger placement on and off the corner of the fret box.

\section{The right amount of force}

When pushing the string down towards the fretboard, you should only use the minimum amount of force necessary to make contact between the string and the fret. Any additional force applied is unnecessary and detrimental to developing good technique. To practice, start with any finger resting on any string at the corner of any fret box. Play the same string in the right hand - it will sound clunky and muted. Continue playing the string in the right hand as you start to push the string down with the left hand finger. Eventually, you will hear the sound change from a clunk to a buzz. Keep pushing down until the sound becomes completely pure, then stop adding pressure and hold the string in place. Take note of how the finger feels now, because you are using just enough pressure to maintain contact between the string and the fret without over-pressing. Rinse and repeat with different fingers, different strings, and different frets. 

\section{The correct finger shape}

%Image of aqueduct showing arched shape

\section{Left hand finger numbers}

Like the right hand, the left hand also has a system of shorthand for referring to each finger. However, unlike the right hand, we are going to use numbers to refer to the left hand fingers. 

\begin{itemize}
	\item Index finger is \textit{1st} finger
	\item Middle finger is \textit{2nd} finger
	\item Ring finger is \textit{3rd} finger
	\item Pinky finger is \textit{4th} finger 
\end{itemize}

In sheet music, if you see a small number next to a note, it means hold that note with the corresponding left hand finger. For example, a small 2 next to a note means hold the note with middle finger. If you see a 0 next to any note, it means that note represents an open string.

\begin{figure}[h]
	\begin{center}
		\large
		\lilypondfile{g_major_scale_finger_numbers.ly}
		\caption{\textit{Order of LH fingers to be used: None, middle, none, index, pinky, none, middle, pinky. The ring finger is un-used in this example.}}
	\end{center}
\end{figure}

\section{First two exercises: Finger Drops and Spider Fingers}

\section{Learning left hand note names}
